
\documentclass[a4paper, 12]{article}


\usepackage{amsmath,amssymb} 
\usepackage[french]{babel}
\usepackage{fancyvrb}
                                                                                                                               
\newtheorem{ex}{Exemple}

\newcommand\set[1]{\ensuremath{\{#1\}}}
\newcommand\acc[1]{$\{$\emph{#1}$\}$}
\newcommand\acco{$\{$}
\newcommand\accf{$\}$}



\begin{document}

\title{Raisonnement qualitatif sur les difficultés dans \emph{Kimind}}
\date{}
\maketitle

\section{Introduction}
Ce document explore la façon de qualifier la difficulté d’une partie de \emph{Kimind} pour un certain joueur, compte tenu des caractéristiques de la partie et du joueur.

\section{Variables qualitatives}

Une difficulté notée $D$ ci-dessous prend un valeur dans $Q = $ \acc{very\_difficult, difficult, medium, easy, very\_easy}.


\subsection{Caractéristiques du joueur}

\subsubsection*{Façon de joueur}

\begin{itemize}
\item \emph{plays\_colours\_order(O)} (où \emph{O} est dans \acc{in\_order, reverse\_order})
\item \emph{respects\_colours},
\item \emph{respects\_sum},
\item \emph{strategy(S)} (où \emph{S} est dans \acc{dichotomic, limits})
\end{itemize}


\subsubsection*{Qualités du joueur}
\begin{itemize}
	\item \emph{rapidity(R)} (\emph{R} est dans \acc{fast,  slow})
	\item \emph{skill\_level(L)} (\emph{L} est dans \acc{strong, weak, medium})
	\item \emph{difficulties\_with\_high\_sum} (booléen)
\end{itemize}

\subsection{Caractéristiques de la partie}
Étant donné une partie caractérisée par $(s, o, \omega)$ (somme, ordre des couleurs et solution)

\begin{itemize}
	\item \emph{high\_sum} (booléen),
	\item \emph{game\_difficulty(D)} (en fonction du nombre de solutions),
	\item \emph{type\_solution(S)} (la solution se prête à une stratégie) (\emph{S} est dans \acc{dichotomic, limits}),
	\item \emph{order\_colours\_solution(O)} (respecte l’ordre des couleurs) (\emph{O} est dans \acc{order, reverse\_order}).
\end{itemize}


\section{Qualifier la difficulté d’une partie pour un joueur}

\subsection{Variables qualitatives}

\begin{itemize}
	\item \emph{high\_sum(D)},
	\item \emph{same\_colours\_order(D)},
	\item \emph{reverse\_colours\_order(D)},
	\item \emph{game\_difficulty(D)}.
\end{itemize}

Pour chaque variable ci-dessous, la valeur $D$ dénotant une difficulté est instanciée par des règles de \emph{bk\_partie\_joueur.pl}

\subsection{Règles}
Nous qualifions la difficulté globale d’une partie (\emph{overall\_difficulty(D)}) pour un joueur donné avec un raisonnement qualitatif à base de règles.

\begin{itemize}
\begin{Verbatim}
overall_difficult(D) <- 
	same_colours_order(D1),
	
\end{Verbatim}
\item 
\

\item Un fort joueur a une difficulté globale améliorée d’un cran :
\begin{Verbatim}
player_game_difficulty (D) <- 
	strong_player,
	overall_difficulty(D1),
	improve_quality(D1, D).
\end{Verbatim}
	
\item Un faible joueur a une difficulté globale dégradée d’un cran :
\begin{Verbatim}
player_game_difficulty (D) <- 
	weak_player,
	overall_difficulty(D1),
	degrade_quality(D1, D).
\end{Verbatim}

\end{itemize}


  
\end{document}